\index{BPatch_basicBlock: : : : : : :!}{

The BPatch_loopTreeNode class provides a tree interface to a collection of instances of class BPatch_basicBlockLoop

contained in a BPatch_flowGraph. The structure of the tree follows the nesting relationship of the loops in a

function's flow graph. Each BPatch_­loopTreeNode contains a pointer to a loop (represented by

BPatch_basicBlockLoop), and a set of sub-loops (represented by other BPatch_loopTreeNode objects). The root

BPatch_­loopTreeNode instance has a null loop member since a function may contain multiple outer loops. The outer

loops are contained in the root instance's vector of children.



{

Each instance of BPatch_loopTreeNode is given a name that indicates its position in the hierarchy of loops. The name

of each root loop takes the form of loop_x, where x is an integer from 1 to n, where n is the number of outer loops in

the function. Each sub-loop has the name of its parent, followed by a .y, where y is 1 to m, where m is the number of

sub-loops under the outer loop. For example, consider the following C function:}



{

void foo() {}



{

int x, y, z, i;}



{

for (x=0; x<10; x++) {}



{

for (y = 0; y<10; y++)}



{

...}



{

for (z = 0; z<10; z++)}



{

...}



{

}}



{

for (i = 0; i<10; i++) {}



{

...}



{

}}



{

}





\bigskip



{

The foo function will have a root BPatch_loopTreeNode, containing a NULL loop entry and two BPatch_loopTreeNode

children representing the functions outer loops. These children would have names loop_1 and loop_2, respectively

representing the x and i loops. loop_2 has no children. loop_1 has two child BPatch_loopTreeNode objects, named

loop_1.1 and loop_1.2, respectively representing the y and z loops.



{

BPatch_basicBlockLoop *loop\index{getSources: : : : : : :!}}



{

A node in the tree that represents a single BPatch_basicBlockLoop instance.}



{

std::vector<BPatch_loopTreeNode *> children\index{getSources: : : : : : :!}}



{

The tree nodes for the loops nested under this loop.}



{

const char *name()\index{getSources: : : : : : :!}}



{

Return a name for this loop that indicates its position in the hierarchy of loops.}



{

bool getCallees(std::vector<BPatch_function *> &v, BPatch_addressSpace *p)}



{

This function fills the vector v with the list of functions that are called by this loop.}



{

const char *getCalleeName(unsigned int i)}



{

This function return the name of the ith function called in the loop's body.}



{

unsigned int numCallees()}



{

Returns the number of callees contained in this loop's body.}



{

BPatch_basicBlockLoop *findLoop(const char *name)}



{

Finds the loop object for the given canonical loop name.}



